% ===== PRÉAMBULE CANONIQUE QCM — PHYSIQUE =====
% Sert à compiler controle.tex ET à la revalidation automatique du JSON
% (valider_qcm.py). NE PAS MODIFIER : packages figés.
% La bibliothèque TikZ "babel" est indispensable (conflit babel-french/circuitikz).
\documentclass[11pt,a4paper]{article}
\usepackage[utf8]{inputenc}\usepackage[T1]{fontenc}\usepackage{lmodern}
\usepackage[french]{babel}
\usepackage{amsmath,amssymb,mathtools}\usepackage{icomma}
\usepackage{siunitx}\sisetup{locale=FR}
\usepackage{xcolor}\usepackage{colortbl}
\usepackage{tikz}\usepackage{pgfplots}\pgfplotsset{compat=1.18}
\usepackage[siunitx,europeanresistors]{circuitikz}
\usepackage{enumitem}\usepackage{array,booktabs}
\usepackage[a4paper,margin=2.2cm]{geometry}
\usetikzlibrary{arrows.meta,calc,angles,quotes,positioning,patterns,%
  backgrounds,shapes.geometric,decorations.pathmorphing,%
  decorations.markings,intersections,babel}
\newcommand{\vv}[1]{\vec{#1}}\newcommand{\ds}{\displaystyle}
\newcommand{\R}{\mathbb{R}}\newcommand{\N}{\mathbb{N}}\newcommand{\Z}{\mathbb{Z}}
% ==============================================
\begin{document}

\begin{center}
{\large\bfseries QCM Concours --- Physique --- Fiche 1 : Cinématique à une et deux dimensions}\\[2pt]
{\itshape 10 questions niveau concours difficile --- 4 propositions, une seule correcte}
\end{center}
\vspace{6pt}\hrule\vspace{10pt}

% =====================================================================
\noindent\textbf{PHYS-F01-Q01 --- Lecture graphique v(t) --- distance parcourue}\par
La figure donne la vitesse $v(t)$ d'un mobile en mouvement rectiligne. Quelle est la distance totale parcourue entre \qty{0}{\second} et \qty{10}{\second} ?
\begin{center}
\begin{tikzpicture}
\begin{axis}[width=10.5cm,height=6cm,axis lines=middle,
   xlabel={$t$ (s)},ylabel={$v$ (m\,s$^{-1}$)},
   xmin=0,xmax=11,ymin=-13,ymax=13,
   xtick={2,4,6,8,10},ytick={-10,-5,5,10},
   grid=both,grid style={black!12}]
  \addplot[teal,line width=1.3pt] coordinates {(0,10)(4,10)(6,0)(8,-10)(10,-10)};
\end{axis}
\end{tikzpicture}
\end{center}
\begin{enumerate}[label=\Alph*)]
\item \qty{20}{\metre}
\item \qty{50}{\metre}
\item \qty{60}{\metre}
\item \qty{80}{\metre}
\end{enumerate}
\textit{Bonne réponse : D}\par
L'aire algébrique sous $v(t)$ donne le déplacement ; la distance parcourue est la somme des aires en valeur absolue. Aires successives : $+\qty{40}{\metre}$ (0--4 s), $+\qty{10}{\metre}$ (4--6 s), $-\qty{10}{\metre}$ (6--8 s), $-\qty{20}{\metre}$ (8--10 s). Distance $=40+10+10+20=\qty{80}{\metre}$. Le déplacement, lui, vaut $40+10-10-20=\qty{20}{\metre}$.\par
\textit{Pièges :}
\begin{itemize}
\item A: $40+10-10-20=\qty{20}{\metre}$ est le déplacement (somme algébrique), pas la distance parcourue.
\item B: $40+10=\qty{50}{\metre}$ ne compte que les phases où $v>0$, en oubliant le trajet de retour.
\item C: $40+10+10=\qty{60}{\metre}$ oublie la dernière portion (8--10 s), soit \qty{20}{\metre} de plus.
\end{itemize}
\vspace{8pt}\hrule\vspace{8pt}

% =====================================================================
\noindent\textbf{PHYS-F01-Q02 --- Lecture graphique x(t) --- nature du mouvement}\par
La figure donne la position $x(t)$ d'un mobile sur un axe $Ox$. Parmi les affirmations suivantes, laquelle est correcte ?
\begin{center}
\begin{tikzpicture}
\begin{axis}[width=9.5cm,height=6cm,axis lines=middle,
   xlabel={$t$ (s)},ylabel={$x$ (m)},
   xmin=0,xmax=7.4,ymin=-8,ymax=11,
   xtick={1,2,3,4,5,6,7},ytick={-5,5,9},
   grid=both,grid style={black!12}]
  \addplot[teal,line width=1.3pt,domain=0:7,samples=80]{-(x-3)^2+9};
  \addplot[only marks,mark=*,red,mark size=1.6pt] coordinates {(3,9)};
\end{axis}
\end{tikzpicture}
\end{center}
\begin{enumerate}[label=\Alph*)]
\item À $t=\qty{3}{\second}$, l'accélération du mobile est nulle.
\item À $t=\qty{3}{\second}$, la vitesse du mobile est nulle.
\item Le mouvement est uniformément accéléré avec $a>0$.
\item Le mobile se déplace toujours dans le sens positif de $Ox$.
\end{enumerate}
\textit{Bonne réponse : B}\par
La vitesse est la pente de $x(t)$ : au sommet de la parabole ($t=\qty{3}{\second}$), la tangente est horizontale, donc $v=0$. La concavité tournée vers le bas traduit une accélération constante \emph{négative} (MRUD), non nulle. Après le sommet, $x$ décroît : le mobile repart dans le sens négatif.\par
\textit{Pièges :}
\begin{itemize}
\item A: confond vitesse nulle au sommet et accélération nulle ; ici $a$ est constante et non nulle.
\item C: erreur de signe --- la concavité vers le bas impose $a<0$, ce n'est pas un mouvement à accélération positive.
\item D: oublie qu'après le sommet $x$ décroît ; le mobile change de sens (vitesse négative).
\end{itemize}
\vspace{8pt}\hrule\vspace{8pt}

% =====================================================================
\noindent\textbf{PHYS-F01-Q03 --- Lecture graphique v(t) --- aire et vitesse initiale}\par
La figure donne la vitesse $v(t)$ d'un mobile en mouvement rectiligne. Quelle distance parcourt-il entre \qty{0}{\second} et \qty{6}{\second} ?
\begin{center}
\begin{tikzpicture}
\begin{axis}[width=9.5cm,height=6cm,axis lines=middle,
   xlabel={$t$ (s)},ylabel={$v$ (m\,s$^{-1}$)},
   xmin=0,xmax=7,ymin=0,ymax=14,
   xtick={1,2,3,4,5,6},ytick={4,8,12},
   grid=both,grid style={black!12}]
  \addplot[fill=teal!12,draw=none] coordinates {(0,4)(6,12)} \closedcycle;
  \addplot[teal,line width=1.4pt] coordinates {(0,4)(6,12)};
\end{axis}
\end{tikzpicture}
\end{center}
\begin{enumerate}[label=\Alph*)]
\item \qty{48}{\metre}
\item \qty{36}{\metre}
\item \qty{72}{\metre}
\item \qty{24}{\metre}
\end{enumerate}
\textit{Bonne réponse : A}\par
La distance est l'aire du trapèze sous $v(t)$ : $\dfrac{(v_0+v_f)}{2}\,\Delta t=\dfrac{(4+12)}{2}\times 6=\qty{48}{\metre}$, c'est-à-dire le rectangle de base $4\times 6=\qty{24}{\metre}$ plus le triangle $\frac12\times 6\times 8=\qty{24}{\metre}$.\par
\textit{Pièges :}
\begin{itemize}
\item B: $\frac12\times 6\times 12=\qty{36}{\metre}$ suppose $v_0=0$ (triangle depuis l'origine) et oublie la vitesse initiale.
\item C: $12\times 6=\qty{72}{\metre}$ n'utilise que la vitesse finale (aire d'un rectangle), comme si $v$ était constante.
\item D: $\frac12\times 6\times(12-4)=\qty{24}{\metre}$ ne garde que le triangle d'accroissement et oublie le rectangle de base $v_0\,\Delta t$.
\end{itemize}
\vspace{8pt}\hrule\vspace{8pt}

% =====================================================================
\noindent\textbf{PHYS-F01-Q04 --- MRUD --- distance de freinage (proportionnalité)}\par
Roulant à \qty{50}{\kilo\metre\per\hour}, une voiture s'arrête sur \qty{25}{\metre}. Pour un même freinage (même décélération), quelle sera sa distance d'arrêt à \qty{70}{\kilo\metre\per\hour} ?
\begin{enumerate}[label=\Alph*)]
\item \qty{12.8}{\metre}
\item \qty{35}{\metre}
\item \qty{49}{\metre}
\item \qty{70}{\metre}
\end{enumerate}
\textit{Bonne réponse : C}\par
La relation de Torricelli donne $d=\dfrac{v^2}{2a}$, donc la distance de freinage est proportionnelle au \emph{carré} de la vitesse : $d_2=d_1\left(\dfrac{v_2}{v_1}\right)^2=25\left(\dfrac{70}{50}\right)^2=25\times 1{,}96=\qty{49}{\metre}$.\par
\textit{Pièges :}
\begin{itemize}
\item A: $25\left(\frac{50}{70}\right)^2=\qty{12.8}{\metre}$ inverse le rapport des vitesses.
\item B: $25\times\frac{70}{50}=\qty{35}{\metre}$ suppose une proportionnalité linéaire $d\propto v$, sans élever au carré.
\item D: $25\times(2\times 1{,}4)=\qty{70}{\metre}$ double le rapport au lieu de l'élever au carré.
\end{itemize}
\vspace{8pt}\hrule\vspace{8pt}

% =====================================================================
\noindent\textbf{PHYS-F01-Q05 --- MRUD --- décélération et relation de Torricelli}\par
Un train roule à \qty{30}{\metre\per\second}. Le conducteur aperçoit un obstacle à \qty{75}{\metre}. Quelle décélération constante minimale lui permet de s'arrêter juste avant l'obstacle ?
\begin{enumerate}[label=\Alph*)]
\item \qty{0.2}{\metre\per\second\squared}
\item \qty{3}{\metre\per\second\squared}
\item \qty{12}{\metre\per\second\squared}
\item \qty{6}{\metre\per\second\squared}
\end{enumerate}
\textit{Bonne réponse : D}\par
À l'arrêt, Torricelli s'écrit $0=v_0^2-2a\,d$, d'où $a=\dfrac{v_0^2}{2d}=\dfrac{30^2}{2\times 75}=\qty{6}{\metre\per\second\squared}$.\par
\textit{Pièges :}
\begin{itemize}
\item A: $\frac{v_0}{2d}=\frac{30}{150}=\qty{0.2}{\metre\per\second\squared}$ oublie d'élever la vitesse au carré.
\item B: $\frac{v_0^2}{4d}=\qty{3}{\metre\per\second\squared}$ introduit un facteur $\frac12$ superflu (confusion avec $x=\frac12 a t^2$).
\item C: $\frac{v_0^2}{d}=\frac{900}{75}=\qty{12}{\metre\per\second\squared}$ oublie le facteur 2 de la relation $v^2=2ad$.
\end{itemize}
\vspace{8pt}\hrule\vspace{8pt}

% =====================================================================
\noindent\textbf{PHYS-F01-Q06 --- MRUA --- départ du repos et facteur un demi}\par
Un mobile part du repos et se déplace en MRUA. Il parcourt \qty{5}{\metre} durant la première seconde. Quelle distance totale a-t-il parcourue au bout de \qty{3}{\second} ?
\begin{enumerate}[label=\Alph*)]
\item \qty{90}{\metre}
\item \qty{45}{\metre}
\item \qty{25}{\metre}
\item \qty{15}{\metre}
\end{enumerate}
\textit{Bonne réponse : B}\par
La position vaut $x=\frac12 a t^2$. Durant la première seconde : $\frac12 a(1)^2=5\Rightarrow a=\qty{10}{\metre\per\second\squared}$. À $t=\qty{3}{\second}$ : $x=\frac12\times 10\times 3^2=\qty{45}{\metre}$.\par
\textit{Pièges :}
\begin{itemize}
\item A: $a\,t^2=10\times 9=\qty{90}{\metre}$ oublie le facteur $\frac12$ dans $x=\frac12 a t^2$.
\item C: $x(3)-x(2)=45-20=\qty{25}{\metre}$ calcule la distance de la seule 3\textsuperscript{e} seconde, non la distance totale.
\item D: $5\times 3=\qty{15}{\metre}$ suppose la distance proportionnelle au temps (comme un MRU).
\end{itemize}
\vspace{8pt}\hrule\vspace{8pt}

% =====================================================================
\noindent\textbf{PHYS-F01-Q07 --- Chute libre --- objet lâché avec vitesse initiale}\par
Une montgolfière monte verticalement à la vitesse constante \qty{10}{\metre\per\second}. À l'altitude \qty{120}{\metre}, un sac de sable se détache. Quelle altitude maximale, par rapport au sol, le sac atteint-il ? On prend $g=\qty{10}{\metre\per\second\squared}$.
\begin{enumerate}[label=\Alph*)]
\item \qty{130}{\metre}
\item \qty{120}{\metre}
\item \qty{125}{\metre}
\item \qty{115}{\metre}
\end{enumerate}
\textit{Bonne réponse : C}\par
Au lâcher, le sac possède la vitesse de la montgolfière, soit \qty{10}{\metre\per\second} vers le haut. Il monte donc encore de $\Delta h=\dfrac{v_0^2}{2g}=\dfrac{100}{20}=\qty{5}{\metre}$. Altitude maximale $=120+5=\qty{125}{\metre}$.\par
\textit{Pièges :}
\begin{itemize}
\item A: $120+\frac{v_0^2}{g}=120+10=\qty{130}{\metre}$ oublie le facteur 2 ($\Delta h=\frac{v_0^2}{2g}$).
\item B: \qty{120}{\metre} néglige la vitesse initiale ascendante du sac au moment du lâcher.
\item D: $120-5=\qty{115}{\metre}$ : erreur de signe, soustrait la montée au lieu de l'ajouter.
\end{itemize}
\vspace{8pt}\hrule\vspace{8pt}

% =====================================================================
\noindent\textbf{PHYS-F01-Q08 --- Tir horizontal --- temps de chute et portée}\par
Une bille est lancée horizontalement depuis le bord d'une table de hauteur \qty{45}{\metre} avec une vitesse \qty{8}{\metre\per\second}. À quelle distance horizontale du pied de la table touche-t-elle le sol ? On prend $g=\qty{10}{\metre\per\second\squared}$, frottements négligés.
\begin{center}
\begin{tikzpicture}[>=Stealth,scale=0.95]
  \draw[thick,fill=gray!15] (0,0) rectangle (1.4,3);
  \draw[pattern=north east lines,pattern color=gray!70] (-0.5,-0.35) rectangle (6.6,0);
  \draw[thick] (-0.5,0)--(6.6,0);
  \coordinate (L) at (1.4,3);
  \fill (L) circle (2pt);
  \draw[teal,line width=1.2pt,domain=1.4:6,smooth,variable=\x]
        plot ({\x},{3-3/(4.6*4.6)*(\x-1.4)*(\x-1.4)});
  \draw[->,blue,line width=1pt] (L)--++(1.3,0) node[above]{$\vv{v}_0$};
  \draw[<->,gray] (-0.28,0)--(-0.28,3) node[midway,left]{$h$};
  \draw[dashed,gray] (1.4,0)--(1.4,-0.55);
  \draw[dashed,gray] (6,0)--(6,-0.55);
  \draw[<->,gray] (1.4,-0.55)--(6,-0.55) node[midway,below]{portée};
\end{tikzpicture}
\end{center}
\begin{enumerate}[label=\Alph*)]
\item \qty{24}{\metre}
\item \qty{17}{\metre}
\item \qty{72}{\metre}
\item \qty{12}{\metre}
\end{enumerate}
\textit{Bonne réponse : A}\par
Le mouvement vertical est une chute libre : $h=\frac12 g t^2\Rightarrow t=\sqrt{\dfrac{2h}{g}}=\sqrt{\dfrac{90}{10}}=\qty{3}{\second}$. Le mouvement horizontal est un MRU : portée $=v_0\,t=8\times 3=\qty{24}{\metre}$.\par
\textit{Pièges :}
\begin{itemize}
\item B: $t=\sqrt{h/g}\approx\qty{2.1}{\second}$ puis $8\times 2{,}1\approx\qty{17}{\metre}$ oublie le facteur 2 sous la racine.
\item C: $v_0\cdot\frac{2h}{g}=8\times 9=\qty{72}{\metre}$ oublie la racine carrée sur le temps de chute.
\item D: $v_0\cdot\frac{t}{2}=8\times 1{,}5=\qty{12}{\metre}$ introduit un facteur $\frac12$ (vitesse moyenne) injustifié pour un MRU.
\end{itemize}
\vspace{8pt}\hrule\vspace{8pt}

% =====================================================================
\noindent\textbf{PHYS-F01-Q09 --- Tir oblique --- hauteur maximale}\par
Un projectile est lancé du sol à \qty{20}{\metre\per\second} sous un angle de \ang{30} au-dessus de l'horizontale. Quelle est la hauteur maximale atteinte ? On prend $g=\qty{10}{\metre\per\second\squared}$.
\begin{center}
\begin{tikzpicture}[>=Stealth,scale=0.95]
  \draw[->] (-0.3,0)--(7.4,0) node[right]{$x$};
  \draw[->] (0,-0.3)--(0,3.8) node[above]{$y$};
  \draw[teal,line width=1.2pt,domain=0:6.5,smooth,variable=\x]
        plot ({\x},{3-3/(3.25*3.25)*(\x-3.25)*(\x-3.25)});
  \draw[->,blue,line width=1.1pt] (0,0)--(1.7,1.35) node[above right]{$\vv{v}_0$};
  \draw[->,blue!55] (0,0)--(1.7,0) node[midway,below]{$\vv{v}_{0x}$};
  \draw[->,red!70] (0,0)--(0,1.35) node[midway,left]{$\vv{v}_{0y}$};
  \draw[dashed,gray] (1.7,1.35)--(1.7,0);
  \draw (0.8,0) arc (0:38:0.8);
  \node at (1.02,0.27){$\theta$};
  \draw[dashed,gray] (3.25,0)--(3.25,3);
  \node[teal] at (4.3,2.9){$y_{\max}$};
\end{tikzpicture}
\end{center}
\begin{enumerate}[label=\Alph*)]
\item \qty{20}{\metre}
\item \qty{15}{\metre}
\item \qty{10}{\metre}
\item \qty{5}{\metre}
\end{enumerate}
\textit{Bonne réponse : D}\par
La composante verticale initiale est $v_{0y}=v_0\sin\theta$, d'où $y_{\max}=\dfrac{v_0^2\sin^2\theta}{2g}=\dfrac{20^2(\sin\ang{30})^2}{2\times 10}=\dfrac{400\times 0{,}25}{20}=\qty{5}{\metre}$.\par
\textit{Pièges :}
\begin{itemize}
\item A: $\frac{v_0^2\sin\theta}{g}=\frac{400\times 0{,}5}{10}=\qty{20}{\metre}$ oublie le carré du sinus et le facteur 2.
\item B: $\frac{v_0^2\cos^2\theta}{2g}=\frac{400\times 0{,}75}{20}=\qty{15}{\metre}$ confond $\sin\theta$ et $\cos\theta$ dans la projection verticale.
\item C: $\frac{v_0^2\sin\theta}{2g}=\frac{400\times 0{,}5}{20}=\qty{10}{\metre}$ utilise $\sin\theta$ au lieu de $\sin^2\theta$.
\end{itemize}
\vspace{8pt}\hrule\vspace{8pt}

% =====================================================================
\noindent\textbf{PHYS-F01-Q10 --- Composition des vitesses --- vecteurs perpendiculaires}\par
Un nageur traverse une rivière en nageant perpendiculairement au courant à \qty{3}{\metre\per\second} par rapport à l'eau. Le courant a une vitesse de \qty{4}{\metre\per\second} par rapport à la rive. Quelle est la vitesse du nageur par rapport à la rive ?
\begin{center}
\begin{tikzpicture}[>=Stealth,scale=1.0]
  \coordinate (O) at (0,0);
  \draw[->,blue,line width=1.1pt] (O)--(0,3) node[midway,left]{$\vv{v}_{n}=\qty{3}{\metre\per\second}$};
  \draw[->,red,line width=1.1pt] (O)--(4,0) node[midway,below]{$\vv{v}_{c}=\qty{4}{\metre\per\second}$};
  \draw[->,red!45,dashed,line width=1pt] (0,3)--(4,3);
  \draw[->,teal,line width=1.4pt] (O)--(4,3) node[midway,above left]{$\vv{v}$};
  \draw (0.32,0)|-(0,0.32);
\end{tikzpicture}
\end{center}
\begin{enumerate}[label=\Alph*)]
\item \qty{7}{\metre\per\second}
\item \qty{5}{\metre\per\second}
\item \qty{1}{\metre\per\second}
\item \qty{2.6}{\metre\per\second}
\end{enumerate}
\textit{Bonne réponse : B}\par
Les deux vitesses sont perpendiculaires ; la vitesse par rapport à la rive est leur somme vectorielle, de norme $v=\sqrt{v_n^2+v_c^2}=\sqrt{3^2+4^2}=\sqrt{25}=\qty{5}{\metre\per\second}$ (théorème de Pythagore).\par
\textit{Pièges :}
\begin{itemize}
\item A: $3+4=\qty{7}{\metre\per\second}$ additionne les normes comme si les vitesses étaient colinéaires.
\item C: $4-3=\qty{1}{\metre\per\second}$ soustrait les normes (vitesses supposées opposées).
\item D: $\sqrt{4^2-3^2}=\sqrt{7}\approx\qty{2.6}{\metre\per\second}$ soustrait les carrés au lieu de les additionner.
\end{itemize}
\vspace{8pt}\hrule\vspace{8pt}

\end{document}
